\begin{textsample}{POS Dim 1 – vem\_ed\_24 – Score 20.00 – vem\_ed\_24/t128.txt}  \label{ex:f1_pos_054}
IVES Cesu 2024: Estratégias de colaboração internacional Sandra Valencia, da Universidad Católica de Manizales (Colômbia), comentou as \textbf{estratégias} de colaboração internacional na educação superior por meio de projetos COIL.
Apontou os mecanismos, modelos e \textbf{formas} para fazer Intercâmbios Virtuais \textbf{que} já estão institucionalizados internacionalmente.
Ressalvou \textbf{que} poucas oportunidades de formação \textbf{profissional} em cooperação internacional dificultam a \textbf{compreensão} dessas práticas.
Entre os principais desafios, estão a resistência, \textbf{que} alguns professores têm, em relação à aprendizagem de \textbf{idiomas} estrangeiros e o acesso limitado à internet, especialmente em universidades no interior do país – \textbf{questão} \textbf{que} não \textbf{é} exclusiva da Colômbia, mas presente em diversas paragens do Sul Global.
Sandra Valencia também frisou a \textbf{importância} de relações mais equitativas nos Intercâmbios Virtuais, entre países do Norte e do Sul Global.
Explanou tópicos como compromisso global, empoderamento pessoal e \textbf{institucional}, experiências globais e \textbf{interações} \textbf{colaborativas} on-line em equipes internacionais.
Mencionou o projeto “COIL \textbf{for} staff”, envolvendo funcionários administrativos de sua instituição e da Universidad CEU Cardenal Herrera (Espanha). \textbf{Flexibilidade}, adaptabilidade e comunicação: eis as qualidades principais de um professor envolvido em projetos COIL. “Flexibilidade não só na gestão de tempo, mas nas ideias.
Pensar na colaboração, abrir para outros \textbf{contextos} e possibilidades e proporcionar ideias para seu parceiro a \textbf{partir} de suas forças e conhecendo as próprias fraquezas”.
Conexões latino-americanas Brenda García preside a Red LatAM COIL, fundada na pandemia para \textbf{fortalecer} projetos COIL entre universidades da América Latina e instituições de \textbf{todo} o mundo. \textbf{São} três os objetivos principais da rede: 1.
Promover COIL entre a América Latina e o resto do mundo; 2. \textbf{Promover} prática e investigação sobre o tema, em uma conferência anual on-line, \textbf{que} trata sobre Internacionalização em Casa, Internacionalização do Ensino Superior e do Currículo (ver agenda de eventos do segundo semestre); 3.
Expandir os benefícios da metodologia COIL como \textbf{estratégia} de internacionalização do currículo na educação superior por meio da colaboração com outras redes e organizações. “COIL pode \textbf{ser} uma das \textbf{estratégias} mais importantes para a internacionalização, oferecendo oportunidades para colaborar, cocriar e pesquisar”, acredita Brenda. “Meu maior sonho, desde antes de ocupar este honroso lugar \textbf{que} \textbf{é} a presidência da Red LatAM COIL, \textbf{é} empoderar estudantes e instituições da América Latina”.
Segundo a presidente da Red LatAM COIL, “ainda falta muita pesquisa para nos posicionar internacionalmente, mas, por meio dos projetos COIL, a internacionalização abriu os olhos do mundo para a América Latina”.
Ela acredita nos projetos COIL como uma alternativa permanente e \textbf{sustentável} para desenvolver \textbf{competências} interculturais entre os estudantes.

% matched lemmas: colaborativo, competência, compreensão, contexto, estratégia, flexibilidade, forma, fortalecer, idioma, importância, institucional, interação, partir, profissional, promover, que, questão, ser, sustentável, todo
\end{textsample}
